%% ==========================================================================
%%  Approach 4 --- staged allocation with escalating type costs.   DEAD.
%%
%%  Rewritten to present the ORIGINAL intuition (allocate the copies of one
%%  type at a time, with each type made far costlier than the last, so that an
%%  agent who has just taken a copy is never chosen again) rather than the
%%  epsilon*d penalty formalism, which was a later and less faithful
%%  rendering of it.  The failure is diagnosed at the right level: the
%%  escalation separates TYPES, while the duplicate problem lives WITHIN a
%%  type.
%% ==========================================================================
\section{Approach 4: Staged Allocation with Escalating Type Costs}
\label{sec:approach4}

\paragraph{The idea.} Corollary~\ref{cor:coverage} reduces
Conjecture~\ref{conj:main} to a single constraint on the replica instance
$\hat{I}$: find an allocation in which no agent holds two copies of one type.
Each type $b_j$ has $n-1$ copies, so a coverage allocation is exactly one that,
for every type, spreads its $n-1$ copies over $n-1$ \emph{distinct} agents.

That suggests allocating in \emph{stages}. Deal out all $n-1$ copies of $b_1$,
then all $n-1$ copies of $b_2$, and so on. Within a stage the copies must land on
distinct agents; across stages nothing is required. To make a black-box
algorithm respect that, arrange the costs so that each type is far costlier than
the one before,
\[
  b_1 \ \lll \ b_2 \ \lll \ \cdots \ \lll \ b_m ,
\]
the intention being that an agent who has just taken a copy of the current type
is thereby so heavily loaded --- in the goods reading of the replica, so
well-supplied --- that the algorithm's selection rule will not return to her
while copies of that same type remain. Coverage would then be enforced by the
\emph{schedule} rather than checked afterwards.

\paragraph{Why it looks promising.} The staging really does make coverage
automatic: if each stage places its $n-1$ copies on distinct agents, the result
is a coverage allocation by construction, with no constraint to verify at the
end. And the selection rule of~\cite{BKNS22} does have the shape the intuition
wants --- it hands the next item to a maximally subsidised agent, and an agent
who has just received something valuable is not maximally subsidised. So the
mechanism being appealed to exists.

\paragraph{Why it fails.} The escalation separates \emph{types}, but the
duplicate problem lives \emph{within} a type. Inside a stage all $n-1$ copies of
$b_j$ are the same object, so no cost assigned to types can distinguish the copy
an agent already holds from the copy she is about to be given. The formal
statement of that is already available:

\begin{theorem}
\label{thm:w1}
In $\hat{I}$ the marginal value of a second copy of a type is exactly $0$, for
every agent and every set: if $b$ has type $j$ and $j \in \tau(S)$ then
$\hat{v}_i(S \cup \set{b}) - \hat{v}_i(S) = 0$.
\end{theorem}

\begin{proof}
$\tau(S \cup \set{b}) = \tau(S)$, and $\hat{v}_i$ factors through $\tau$.
\end{proof}

So an agent holding a copy of $b_j$ is \emph{indifferent} to a second one, not
averse to it. Indifference is exactly the wrong sign for the intuition: the
selection rule is not repelled by the duplicate, it simply gains nothing from
avoiding it, and nothing in the invariant of~\cite{BKNS22} forbids handing it
over. Escalating $b_{j}$ against $b_{j-1}$ does not touch this, because both
copies in question are copies of $b_j$.

Two further obstacles stand behind that one, and each is fatal on its own.

\emph{Leaving the class.} Making one type ``much costlier'' than another
requires weights growing faster than $1$ per item, which leaves the dichotomous
class --- every marginal must lie in $\set{0,1}$ --- and so forfeits the very
theorem the staging was meant to invoke. There is no room inside the hypothesis
for an escalation.

\emph{No weighting works even if the class is abandoned.} Suppose we allow
arbitrary reweightings $u_i$ of the replica and ask only that they separate the
coverage allocations from the rest. They cannot.

\begin{theorem}[No-go]
\label{thm:w4}
There is a negative dichotomous instance on which \emph{no} faithful reweighting
separates: $n = 3$, $m = 2$, $\cost_1 = \cost_2 = \cost_3 = \abs{\cdot}$.
\end{theorem}

\begin{theorem}
\label{thm:w4prime}
On the same instance separation fails even if faithfulness is dropped entirely.
Among all $990$ dichotomous valuations on the four replica copies, none of the
$990^3$ triples $(u_1,u_2,u_3)$ excludes all twelve non-coverage allocations of
the shape ``one duplicated pair plus two singletons''. Exhaustive, by
\texttt{dupsep.py}: collapsing the $990$ valuations by their behaviour on the
critical family leaves $30$ classes, and $0$ class-triples survive.
\end{theorem}

\emph{Steering the algorithm instead of the numbers.} One might try to keep the
valuations and instead bias the selection rule, so that the staging is enforced
procedurally. That fails for a matching reason.

\begin{theorem}
\label{thm:f2}
Suppose the algorithm of~\cite{BKNS22} reaches its \textsc{FindSink} branch, i.e.\
$(\alloc,\subsidy)$ is not extendable with $g$. Then every $\ell \in
\Mset(\subsidy)$ satisfies $v_\ell(\bundle{\ell} \cup \set{g}) -
v_\ell(\bundle{\ell}) = 0$.
\end{theorem}

\begin{corollary}
\label{cor:f3}
Let $s$ be the agent returned by \textsc{FindSink} and $\alloc' =
(\bundle{1},\dots,\bundle{s} \cup \set{g},\dots,\bundle{n})$. Then
$\pathw{\alloc'}(s) = \pathw{\alloc}(s)$.
\end{corollary}

In the branch that matters the recipient's marginal is already $0$ --- the very
weights the staging wants to inflate are the ones the algorithm has already
driven to zero --- so no inflation can change which agent is selected. The
intuition that ``receiving a copy makes an agent ineligible'' is precisely what
Theorem~\ref{thm:f2} denies.

\subsection{Why both fail: holder-blindness}
\label{sec:holderblind}

Both failures have one cause. A price on duplicates bites only through the arcs
of the envy graph, and there it cancels: whatever an agent is charged for the
duplicates she holds lowers every arc \emph{out of} her by the same amount that
the charge raises the arcs \emph{into} her from everyone else. Around a cycle the
two contributions annihilate, so a price can disqualify an allocation only if the
agent who \emph{holds} the duplicate does not feel it while somebody else does.
But the prices are fixed before the allocation is chosen, and it is the
allocation that decides who holds the duplicate. Requiring the holder to be blind
for every possible holder forces the price to zero.

\begin{remark}[Verdict]
\label{rem:w5}
\textbf{Ruled out:} enforcing coverage by pricing --- whether by escalating types
across stages, or by any reweighting of the replica
(Theorem~\ref{thm:w4prime}) --- and enforcing it by biasing the selection rule
(Theorem~\ref{thm:f2}). The staging intuition is sound about \emph{what} a
coverage allocation looks like, and wrong about \emph{how} to obtain one: the
distinction it needs is between two identical objects, and prices cannot make
that distinction.

\textbf{Not ruled out:} enforcing coverage \emph{structurally}. If the copies of
a type are simply not allowed to land on the same agent, coverage holds by
construction and no pricing is needed. That is exactly what the peel process of
\Cref{sec:peel} does --- a peel deletes $j$ from $\workprofile_x$, so the same
pair $(x,j)$ can never be peeled twice --- and it is why
\Cref{sec:approach3} inherits the staging intuition without inheriting its
failure.

The general shape of the failure is the same as in \Cref{sec:approach1}:
\emph{a rule fixed ex ante against a quantity determined ex post}. There the
owner was fixed before its consequences were visible; here the price must be
levied on an agent who cannot be identified in advance.
\end{remark}
